\subsection*{转动惯性与角动量概念的建立}

质点若不受合外力，将保持匀速直线运动；转动物体若不受合外力矩，也会保持原有的转动状态。
这两句话分别对应线动量与角动量。角动量并不是把“动量”简单换成“角速度”，因为同一个质点
相对于不同参考点的转动趋势不同。必须同时给出参考点 $O$、质点的位置矢量 $\vect r$ 和动量
$\vect p$，定义
\[
  \vect L_O=\vect r\times\vect p.
\]
它的大小等于 $p$ 乘以参考点到动量作用线的垂距，方向由右手定则确定。因此同一条直线运动，
如果参考点取在运动直线上，角动量为零；参考点离开运动直线，角动量一般不为零。这个事实是
判断角动量方向和选取矩心时最基本的几何依据。

外积与内积的物理作用不同。内积把两个矢量映射成标量，适合描述功、功率和投影；外积把两个
矢量映射成垂直于二者张成平面的轴矢量，适合描述有方向的转动趋势。外积没有交换律，而满足
\[
  \vect a\times\vect b=-\vect b\times\vect a,
  \qquad
  \vect a\times(\vect b+\vect c)=
  \vect a\times\vect b+\vect a\times\vect c.
\]
在直角坐标系中展开行列式，可以同时得到三个分量；作图判断时则应先确定从第一个矢量转向
第二个矢量的最小转角，再用右手螺旋确定正负。只看图形“朝上或朝下”而不检查乘法次序，最
容易把出纸面和入纸面画反。

\subsection*{从 Newton 定律到角动量定理}

对固定在惯性系中的参考点 $O$，有
\[
  \frac{\dd\vect L_O}{\dd t}
  =\frac{\dd\vect r}{\dd t}\times m\vect v
   +\vect r\times\frac{\dd\vect p}{\dd t}.
\]
第一项中两个矢量平行，故为零；利用 Newton 第二定律便得到
\[
  \frac{\dd\vect L_O}{\dd t}=\vect r\times\vect F=\vect M_O.
\]
这就是角动量定理。它与动量定理的对应关系很清楚：合外力决定动量的变化率，合外力矩决定
角动量的变化率。力矩为零并不要求力为零，只要力的作用线通过参考点即可。有心力始终指向
力心，所以相对于力心的力矩恒为零；这也是行星轨道、Coulomb 中心力散射和许多绳约束问题中
角动量守恒的共同原因。

参考点若移动，不能机械套用固定点公式。设参考点位置为 $\vect R(t)$，则相对位置为
$\vect r'=\vect r-\vect R$，求导会出现与参考点速度、加速度有关的附加项。竞赛计算中通常取
固定惯性点、系统质心或固定铰点。选质心的好处是内部相对运动与整体平动能够分离；选固定
铰点的好处是铰链冲量和铰链力矩常可直接消去。

\subsection*{有心力轨道、回归点与面积速度}

质点在有心力作用下运动时，角动量方向固定，所以轨迹位于一个平面内。用极坐标描述，有
\[
  L=mr^2\dot\theta,
  \qquad
  \frac{\dd S}{\dd t}=\frac12r^2\dot\theta=\frac{L}{2m}.
\]
因此面积速度恒定。它不是只对椭圆轨道成立，而是所有中心力运动的直接结论；Kepler 第二
定律只是这一结论在万有引力轨道上的具体表现。

质点到力心距离取极值的位置称为回归点。此时 $\dot r=0$，速度只有横向分量，故
$L=mrv$。若两个回归点距离分别为 $r_1,r_2$，角动量守恒给出
\[
  r_1v_1=r_2v_2.
\]
若力还是保守力，再写机械能守恒
\[
  E=\frac12mv^2+U(r),
\]
便可同时求出回归点距离和速度。只写能量守恒通常只能确定速度大小，不能给出横向运动与半径
之间的联系；只写角动量守恒又缺少径向势能信息。将二者合并，等价于研究有效势能
\[
  U_{\rm eff}(r)=U(r)+\frac{L^2}{2mr^2}.
\]
回归点是水平能量线与有效势能曲线的交点，圆轨道还要求交点同时满足
$\dd U_{\rm eff}/\dd r=0$。这一图像能够区分束缚轨道、散射轨道与不存在回归点的单调运动。

对平方反比引力 $U(r)=-GMm/r$，消去回归点速度可得到关于 $1/r$ 的二次方程。二次方程的
两个根对应近日点与远日点；椭圆的力心必须位于焦点而不是几何中心。近日点距离较小，速度较
大；远日点距离较大，速度较小。图中若把太阳放在椭圆中心，虽然仍能画出闭合曲线，却已经失去
了这一动力学含义。

\subsection*{质点系、内力矩与刚体自转}

对 $N$ 个质点，系统关于 $O$ 的角动量为
\[
  \vect L_O=\sum_{i=1}^{N}\vect r_i\times m_i\vect v_i.
\]
把 $\vect r_i=\vect R_C+\vect r'_i$、$\vect v_i=\vect V_C+\vect v'_i$ 代入，并利用
$\sum m_i\vect r'_i=0$、$\sum m_i\vect v'_i=0$，得到
\[
  \vect L_O=\vect R_C\times M\vect V_C+\vect L_C.
\]
第一项描述质心的轨道角动量，第二项描述系统相对质心的内部角动量。这个分解对可形变质点系
也成立；只有进一步写成 $\vect L_C=I\vect\omega$ 时，才需要刚体、定轴或主轴等附加条件。

系统内部两质点之间的作用力若满足 Newton 第三定律并沿两点连线，则一对内力对任意参考点的
总力矩为零。因此
\[
  \frac{\dd\vect L_O}{\dd t}=\sum\vect M_O^{\rm ext}.
\]
若外力矩为零，系统总角动量守恒。碰撞过程中内力很大，却成对出现；只要外力矩在短暂碰撞
时间内的冲量矩可以忽略，碰撞前后仍可使用角动量守恒。这正是哑铃碰撞、打击与带铰链碰撞中
常选碰撞点或铰点为矩心的理由。

刚体绕定轴转动时，每个质元的速度为 $v_i=\omega r_{i\perp}$，于是
\[
  L=\sum m_i r_{i\perp}^2\omega=I\omega,
  \qquad I=\int r_\perp^2\,\dd m.
\]
转动惯量不仅取决于质量和形状，还取决于转轴。平行轴定理只能在两根平行轴之间使用；挖孔
问题可以把被挖去部分视为负质量，再与完整物体的转动惯量相加。

\subsection*{角动量守恒与陀螺进动}

高速自转陀螺的自转角动量近似沿对称轴。重力矩与角动量近似垂直时，主要改变角动量的方向，
而不是它的大小。设支点到质心距离为 $r$，自转角速度为 $\omega_0$，绕对称轴转动惯量为
$I_0$，稳定进动角速度为 $\omega_p$，则
\[
  Mgr\sin\theta=(I_0\omega_0)\omega_p\sin\theta,
  \qquad
  \omega_p=\frac{Mgr}{I_0\omega_0}.
\]
力矩方向必须垂直于由 $\vect r$ 与 $M\vect g$ 张成的平面；在二维投影中应画成出、入纸面
符号，而不能随意画成平面内箭头。$\vect L$ 的端点绕竖直轴运动，因而扫出圆锥面。上述公式
依赖高速自转和稳定进动近似；若自转不够快，还会出现明显章动，不能把轴倾角视为常量。

\subsection*{叉乘计算与方向的双重检查}

三维角动量图最容易错在“纸面内位置正确，法向方向却反了”。一种可靠办法是先用右手定则判断，
再用分量计算复核。若纸面取为 $xy$ 平面，向右、向上分别为 $+x,+y$，则
\[
 \vect r=(x,y,0),\qquad \vect p=(p_x,p_y,0),
\]
从而
\[
 L_z=xp_y-yp_x.
\]
$L_z>0$ 表示出纸面，$L_z<0$ 表示入纸面。图中若 $\vect r$ 指向右上而 $\vect p$ 水平向右，
便有 $p_y=0$、$y>0$，所以 $L_z=-yp_x<0$，必须画成叉号而不是圆点。力矩
$\vect\tau=\vect r\times\vect F$ 也用同一规则。凡图示给出了三者中的两个方向，第三个方向
就不再是排版选择，而是由叉乘唯一决定。

角动量的单位为 $\mathrm{kg\,m^2\,s^{-1}}$，力矩的单位为
$\mathrm{N\,m}=\mathrm{kg\,m^2\,s^{-2}}$，因此 $\Delta\vect L=\int\vect\tau\,\dd t$
在量纲上自然成立。短时冲击中，不能因为碰撞力很大就武断地说角动量守恒；真正需要比较的是
外力矩的冲量矩与系统角动量改变量。若以碰撞点为矩心，接触冲量对该点力臂为零，往往能得到
最简方程；换到其他矩心后，同一个接触冲量通常不能消去。

\subsection*{中心力轨道的进一步判据}

在极坐标中，单位质量角动量 $h=r^2\dot\theta$ 为常量。令 $u=1/r$，中心力轨道可写成 Binet
形式
\[
 \frac{\dd^2u}{\dd\theta^2}+u=-\frac{F(1/u)}{mh^2u^2}.
\]
它把关于时间的二维运动转化为关于方位角的一维方程。对于万有引力
$F=-GMm/r^2$，右端为常量，解为
\[
 r=\frac{p}{1+e\cos(\theta-\theta_0)}.
\]
$0\le e<1$ 为椭圆，$e=1$ 为抛物线，$e>1$ 为双曲线；圆轨道只是 $e=0$ 的特殊椭圆。若图画
椭圆轨道，力心必须位于焦点而不是凭视觉放在“偏离中心的任意位置”。椭圆半长轴为 $a$、半短
轴为 $b$ 时，焦距满足 $c=\sqrt{a^2-b^2}$，可直接作为作图检查式。

有效势
\[
 U_{\rm eff}(r)=U(r)+\frac{L^2}{2mr^2}
\]
把角向运动的影响表现为离心势垒。径向速度为零的两个位置是近拱点和远拱点，它们由
$E=U_{\rm eff}$ 决定；这两个点并不意味着合力为零，只意味着当时半径停止增大或减小。稳定圆
轨道还需同时满足 $U'_{\rm eff}(r_0)=0$ 与 $U''_{\rm eff}(r_0)>0$。这样便把轨道图、能量图和
角动量守恒连成同一套判断方法。

\subsection*{开普勒定律的角动量解释}

行星在太阳引力下运动时，力始终沿日心连线，故对太阳的力矩为零。面积速度
\[
 \frac{\dd A}{\dd t}=\frac12r^2\dot\theta=\frac{L}{2m}
\]
保持不变，这就是“相等时间扫过相等面积”。行星靠近近日点时半径减小，为维持同一面积速度，
角速度和线速度都必须增大；远日点则相反。图中太阳应放在椭圆焦点，近日点、远日点应位于长轴
两端，扫面积扇形的两条边必须都从太阳出发。

椭圆轨道的总能量
\[
 E=-\frac{GMm}{2a}
\]
只由半长轴 $a$ 决定。结合角动量关系可得轨道周期
$T^2=4\pi^2a^3/(GM)$，说明同一中心天体周围所有椭圆轨道都满足 $T^2\propto a^3$。这里的
$a$ 是半长轴，不是近日点距离，也不是焦点到椭圆中心的距离。把这几个几何量分清，才能正确地
把轨道图与公式对应起来。

角动量守恒还给出一个常用的极限判断：若粒子沿过矩心的直线运动，则
$L=mvb$，其中 $b$ 是矩心到运动直线的垂直距离，而不是矩心到粒子的瞬时距离。粒子虽在移动，
只要直线和速度不变，$b$ 不变，角动量也不变。这一结果常用于散射、掠过碰撞和冲量题，可避免
每个时刻都重新计算叉乘。

最后要始终写明角动量所取的参考点。同一粒子对不同点的角动量一般不同；若矩心本身加速，直接
套用 $\dd\vect L/\dd t=\sum\vect M$ 还会出现附加项。竞赛题优先选择惯性系中固定点、质心或
碰撞作用线经过的点，正是为了让定理形式最简并消去未知力矩。

\subsection*{参考点改变时的角动量}

设 $O$、$O'$ 为同一惯性系中的两个固定点，$\vect a=\overrightarrow{OO'}$，系统总动量为
$\vect P$，则
\[
 \vect L_{O'}=\sum_i(\vect r_i-\vect a)\times m_i\vect v_i
 =\vect L_O-\vect a\times\vect P.
\]
所以“对某点角动量为零”不能不加说明地推广到另一点。若系统总动量为零，角动量才与固定参考点
的选择无关。对质心 $C$ 使用
\[
 \vect L_O=\vect R_C\times M\vect V_C+\vect L_C
\]
可以把质心轨道运动和内部运动分开，这在爆炸、碰撞和变形系统中尤其有用。

力矩也有相应变换
\[
 \vect M_{O'}=\vect M_O-\vect a\times\vect F_{\rm ext}.
\]
当合外力为零时，外力矩对所有固定参考点相同；当合外力不为零时，换矩心会改变外力矩。很多题
选择支点、碰撞点或力的作用线上的点，是因为未知力对该点力矩为零，而不是因为角动量定理对任意
运动点都能保持最简单形式。

三维刚体中更一般的关系是 $\vect L_C=\mathbf I_C\boldsymbol\omega$，其中
$\mathbf I_C$ 是转动惯量张量。只有绕主轴转动时，$\vect L_C$ 才与
$\boldsymbol\omega$ 平行；对任意轴，二者可能有夹角，并使刚体出现进动或章动。定轴问题中写
$L=I\omega$，实质上已经使用了轴方向固定并只取该方向分量的条件。
